\chapter{TECHNICAL CHANGES FROM REVISION 7}
\label{app:errata}

This appendix records the specific technical differences between this revision
and revision~7, with the equation numbers of the earlier document given for
reference.

\section{CHAPTER 3: THE DMSO STABILITY PROOF}

\paragraph{B.1.1 Bound on $\norm{I-\Gamma\phi\phi^T}$ (was Eq.~22).}
Revision~7 used
$\norm{(I-\Gamma\phi\phi^T)\Wt}^2\le(1-\Gamma\norm{\phi}^2)^2\norm{\Wt}^2$. The
matrix $I-\Gamma\phi\phi^T$ is symmetric with eigenvalue $1$ of multiplicity
$l-1$ and eigenvalue $1-\Gamma\norm{\phi}^2$ of multiplicity one, so its induced
$2$-norm is $\max\{1,\abs{1-\Gamma\norm{\phi}^2}\}$, which equals $1$ whenever
$\Gamma\norm{\phi}^2\le2$. The factor used was strictly smaller than the true
norm, so the step was not conservative but invalid. Substituting the correct
value makes the coefficient of $\norm{\Wt_k}^2$ in the former Eq.~(23) equal to
$2/\alpha+3\phi_{\max}^2>0$, and the argument cannot close in that form.

\paragraph{B.1.2 Gain conditions (was Eqs.~26--27).}
Substituting the choice $\alpha=1/(3\phi_{\max}^2)$ into the former Eq.~(26)
gives
\begin{equation}\label{eq:persq}
\frac{2}{\alpha}-6\frac{\Gamma}{\alpha}\phi_{\max}^2
+3\frac{\Gamma^2}{\alpha}\phi_{\max}^4+3\phi_{\max}^2
=9\phi_{\max}^2\bigl(1-\Gamma\phi_{\max}^2\bigr)^2\ \ge\ 0 ,
\end{equation}
a perfect square, whereas the condition required the expression to be
nonpositive. The only admissible point is the boundary
$\Gamma=1/\phi_{\max}^2$, at which the coefficient is zero rather than negative.
The reported condition $\Gamma\le1/\phi_{\max}^2$ was therefore satisfied by no
interior point and excluded small gains, which is the opposite of the correct
behavior established in Theorem~\ref{thm:main} and Corollary~\ref{cor:theta1}.

\paragraph{B.1.3 Weight bound (was Eq.~33).}
The bound inherited the denominator
$-9\phi_{\max}^2(1-\Gamma\phi_{\max}^2)^2\le0$ from \eqref{eq:persq} and was
therefore either undefined or infinite. It is replaced by the honest statement
of Theorem~\ref{thm:main}, which bounds $\Wt_k^T\phi_k$, together with
Section~\ref{sec:weightbound}, which obtains boundedness of $\Wt_k$ from
persistency of excitation or $\sigma$-modification.

\paragraph{B.1.4 Matrix norm (was above Eq.~23).}
Revision~7 used $\norm{A}\le\lmax(A)$. The induced $2$-norm is the largest
singular value, and $\sigmax(A)\ge\rho(A)$, with equality for normal $A$ (and
for some non-normal $A$).
For $A=F-\Km$ the inequality is generally strict and $\lmax(A)$ may be complex.
See Remark~\ref{rem:nonnormal}.

\paragraph{B.1.5 The matrix $Q$ (was after Eq.~28).}
With $P=I$ in the discrete Lyapunov equation $A^TPA-P=-Q$, one has
$Q=I-A^TA$ and $\lmin(Q)=1-\sigmax(A)^2$, so $Q$ is determined by $A$ and is not
independently selectable. The former condition then reduces to
$\sigmax(A)^2\le1/(3+9\Gamma^2\phi_{\max}^4)$, requiring
$\sigmax(A)\lesssim0.29$ at $\Gamma\phi_{\max}^2=1$. The corrected condition is
$\sigmax(A)<1$.

\paragraph{B.1.6 Structural cause.}
Applying Young's inequality to the first difference before evaluating the trace
difference discards the cancellations established in
Proposition~\ref{prop:exact}, which are precisely what the adaptation law was
constructed to produce.

\paragraph{B.1.7 What the Lyapunov argument concludes.} \rev{The first
difference is negative outside a set that bounds $e_k$ and $\psi_k$ but not the
components of $\Wt_k$ orthogonal to $\phi_k$, so that set is not compact in the
observer state $(e_k,\Wt_k)$ and the Lyapunov extension theorem does not apply to
it directly. Theorem~\ref{thm:main} now states a mean-square bound
\eqref{eq:msbound} that needs no bound on the weights, and uniform ultimate
boundedness under a bound on $\Wt_k$ from Section~\ref{sec:weightbound}. The
same applies to Corollary~\ref{cor:deadbeat}. Assumption~\ref{as:nn} now states
that the plant trajectory remains in the compact set on which the approximation
holds, and Remark~\ref{rem:nonnormal} states that a weighted Lyapunov function
requires a correspondingly weighted adaptation law.}

\section{CHAPTER 6: THE EXTRA CONTROL DESIGN}

\paragraph{B.2.1 Missing coupling term (was Eqs.~97 and 100).}
The former Eq.~(97) was obtained by substituting $x_3=\bar x_3$ from the former
Eq.~(93), and was then reused in Eq.~(100) in a setting where $x_3\ne\bar x_3$
by construction. Carrying $x_3=\bar x_3+\bar e_3$ through leaves a term
$+A_2\bar e_3$ in $\dot{\bar e}_2$ that did not appear in the error dynamics or
in the Lyapunov derivative. In the present design this term is retained as
$A_2z_3$ in \eqref{eq:z2dot} and cancelled by the $-A_2z_2$ term in
\eqref{eq:alpha3}.

\paragraph{B.2.2 Undamped error coordinates (was Eqs.~122 and 138--140).}
In revision~7 the cross terms $\bar e_1^T\bar e_2$ and $\bar e_3^T\bar e_4$
cancelled, removing $\bar e_1$ and $\bar e_3$ from the Lyapunov derivative
entirely, and the stated conditions involved only $\norm{\bar e_2}$ and
$\norm{\bar e_4}$. Uniform ultimate boundedness requires the derivative to be
negative outside a compact set in the full error space, which cannot hold in a
direction absent from the derivative. Since $\dot{\bar e}_1=\bar e_2$ is a pure
integrator, a bounded but nonvanishing $\bar e_2$ permits $\bar e_1$ to ramp.
The four-step design of Section~\ref{sec:cffb} assigns damping to all four
coordinates.

\paragraph{B.2.3 Circular approximation targets (was Eqs.~96 and 104).}
The former $F_1$ contained $B_1u_e$, with $u_e$ a function of $\hat F_2$, while
the former $F_2$ contained $-\dot{\bar x}_4$, which contains derivatives of
$\hat F_1$ and hence of the weight estimates. Neither target was a fixed
continuous function of its network input, so the universal approximation
hypothesis did not apply. The targets \eqref{eq:F1} and \eqref{eq:F2} contain no
control signal and no adaptive parameter; the command filters of
Section~\ref{sec:cmdfilter} supply the required derivative signals as filter
states and break the algebraic loop, as detailed in Remark~\ref{rem:order}.

\paragraph{B.2.4 Unmatched first-layer weights (was Eqs.~122 and 134).}
The hidden-layer activation error $\tilde\sigma_{i2}$ was bounded by a constant
$W_{i2m}\sigma_{i2m}$, so the first-layer weight error never appeared in the
Lyapunov derivative and the first-layer update law was unjustified by the
analysis. Equation \eqref{eq:nnerr} uses the standard two-layer Taylor form in
which both layers' weight errors appear linearly and are matched term by term by
the update laws \eqref{eq:upd2}--\eqref{eq:upd1}.

\paragraph{B.2.5 Sign convention (was Eqs.~111--114).}
The two layers carried opposite overall signs in the update laws. Both laws in
\eqref{eq:upd2}--\eqref{eq:upd1} now use the same convention.

\paragraph{B.2.6 Riccati equation (was Eq.~83).}
The trailing factor was written as $A$, a carryover from continuous-time
notation. It is $F$, the discretized system matrix; see \eqref{eq:dare}.

\paragraph{B.2.7 Gain conditions and network residual.} \rev{The state $z_2$
carries two cross terms in the Lyapunov derivative, the filter error and the
residual of network one, so Theorem~\ref{thm:control} requires $c_2>1$, as its
proof always concluded. The Taylor residual $w_i$ of \eqref{eq:nnerr} contains
the weight errors and is bounded by $w_{i0}+w_{i1}\norm{\Wt_i}_F$, not by a
constant; the $\sigma$-modification leakage absorbs the weight-dependent part
when $\kappa_{ij}>2w_{i1}^2$, which is now part of \eqref{eq:gaincond}.}

\section{FRAMING AND REPORTING}

\paragraph{B.3.1 Observer or filter.}
With $y_k=x_k$ the full state is measured. The method is a filter and
uncertainty identifier rather than an observer in the unmeasured-state sense,
and this is now stated explicitly.

\paragraph{B.3.2 Kalman filter comparison.}
Comparing against a Kalman filter that is given no mechanism to accommodate a
deterministic modeling error establishes little. Three baselines are now used,
including an augmented-state filter that estimates the uncertainty as a random
walk, and improvement claims are stated against that one.

\paragraph{B.3.3 Aggregated percentages.}
An average percentage improvement taken over heterogeneous cases is not a
meaningful statistic. Metrics are now reported per case.

\paragraph{B.3.4 Complementary filter result.}
The finding that the complementary filter outperforms both model-based
estimators when those are fed raw accelerometer tilt is retained and explained:
it is a statement about measurement construction, not about estimator quality.
\rev{Re-examined on the hardware data in Section~\ref{sec:implresults}, it holds
for the DMSO but not for a Kalman filter whose measurement variance reflects
the actual accelerometer tilt error.}

\section{MODEL AND RESULTS}

\paragraph{B.4.1 Plant model.} \rev{Four corrections enter the simulations of
Chapter~\ref{ch:results}. The motor torque reacts on the pendulum with positive
sign in the counterclockwise moment balance \eqref{eq:pendm}; revision~7 had the
opposite sign, which changes the back-EMF and input terms of \eqref{eq:nl1},
\eqref{eq:lin1}, and \eqref{eq:linss}. The corrected input vector satisfies the
angular-momentum check of Remark~\ref{rem:momentum}, and the free swing of
Section~\ref{sec:modelsim} is lightly rather than strongly damped. The remaining
three are in Section~\ref{sec:simmodel}: the back-EMF acts on the wheel rate
relative to the body, $\dot x/r+\dot\phi$; $k_e=\SI{0.307}{V.s/rad}$ from the
free-run data in place of \SI{0.00361}{V.s/rad}; and $I_p$ is taken about the
center of gravity. The sensor model now includes the accelerometer kinematics
and the relative encoder measurement $x+r\phi$. Two smaller items: revision~7
stated that $A_4=0$ in the nominal model, but $A_4$ is the gravity coefficient
of the tilt equation; and Table~\ref{tab:phys} lists
$I_w=\SI{1.441e-4}{kg.m^2}$ while the simulations, here and in revision~7's
programs, use $M_wr^2/2=\SI{1.183e-4}{kg.m^2}$.}

\paragraph{B.4.2 Parameter perturbation.} \rev{The torque-constant error of the
uncertainty cases is reduced from $\times2.5$ to $\times1.25$, and $l$ is left
unperturbed, as in the revision~7 extra-control plant. The revision~7
perturbation is still stabilized by every controller on the corrected plant,
but its input-gain error makes the observer bound vacuous during the initial
recovery ($\epsilon_N=0.42$ against $0.013$); see Section~\ref{sec:setup}.}

\paragraph{B.4.3 Command-filtered design.} \rev{The design of
Section~\ref{sec:cffb} is not realizable for the TWIP as written; see
Remark~\ref{rem:realizability} and Section~\ref{sec:cfinstab}. It is corrected in
Section~\ref{sec:flatcf} by three changes: tracking the flat output
$y=x_1-b\,x_3$; using $\eta=x_2-b\,x_4$, whose derivative contains no control, as
the second coordinate, so that $\alpha_2$ no longer contains $B_1u$ and the
virtual-control gain is $a_2=A_2-bA_4$ rather than $A_2$; and weighting the
attitude coordinates by $w=a_2^2$ in the Lyapunov function. The applied
control is independent of $u_{\mathrm{nom}}$ in both forms, since it cancels in
\eqref{eq:ue}.}

\paragraph{B.4.4 Noise table.} \rev{The units of Table~\ref{tab:noise} are
corrected to variances, the tilt white-noise variance is the value used in
simulation ($4.0474\times10^{-6}$ rather than $4.0397\times10^{-6}$~rad$^2$),
and the walking-bias entries are identified as steady-state variances.}
